% session_Q1_results.tex -- Companion note for PCF-2 v1.2 absorption.
% Empirical evidence for Conjecture B4 at d=4 and op:xi0-degree-d at d=4.
%
% Standalone, drop-in section for PCF-2 v1.2.

\documentclass[11pt]{amsart}
\usepackage{amsmath, amssymb, amsthm}
\usepackage{geometry}
\geometry{margin=1in}
\usepackage[hidelinks]{hyperref}
\usepackage{cleveref}

\newtheorem{theorem}{Theorem}
\newtheorem{conjecture}[theorem]{Conjecture}
\newtheorem{proposition}[theorem]{Proposition}
\newtheorem{remark}[theorem]{Remark}
\newtheorem{definition}[theorem]{Definition}

\title{Empirical evidence for the sharp Conjecture~B4 at $d=4$ and the
universal $\xi_0$-identity $\xi_0(b)=d/\beta_d^{1/d}$}
\author{papanokechi (with the SIARC research pipeline)}
\date{2026-05-01}

\begin{document}
\maketitle

\begin{abstract}
We extend the PCF-2 v1.1 quartic test of Conjecture~B4 (PCF-2 v1.1,
Sessions~A2/B/C1) from $d=3$ to $d=4$ on a lex-window catalogue of
$60$ Z-primitive irreducible quartics. The empirical
$A$-distribution is unimodal at $A=8$: $60/60$ families lie within
$0.10$ of $A=2d=8$, mean $\overline{A}=7.954$, $\sigma_A=0.0037$.
We additionally derive and verify a universal Newton-polygon
characteristic root $\xi_0(b)=d/\beta_d^{1/d}$ at $d=4$, recovering
the proven $d=2$ identity $\xi_0=2/\sqrt{\beta_2}$ and supplying
\emph{op:xi0-degree-d} from \textsc{channel-theory-v11}. Together, the
two results refine \emph{op:d2-anomaly}: the $\{2d-1,2d\}$
splitting at $d=2$ is so far the unique multi-branch instance, with
$d\in\{3,4\}$ both populating only the sharp branch $A=2d$ within the
lex-window catalogue.
\end{abstract}

\section{Setup and conjectures}

Recall PCF-2 v1.1, Conjecture~B4 (sharp form, Session~B):
\begin{conjecture}[B4, sharp]\label{conj:B4-sharp}
For every Z-primitive irreducible $b(n)\in\mathbb{Z}[n]$ of degree
$d\ge 3$ with $a_n=1$, the WKB-decay exponent of the partial-tail
errors $\delta_n=L_n - L$ satisfies $A=2d$ in the asymptotic ansatz
$\log|\delta_n|\sim -A\,n\log n + \alpha n - \beta\log n + \gamma$.
\end{conjecture}
At $d=2$ Conjecture~B4 fails as stated: PCF-1 v1.3 Theorem~5 proves
the split $A\in\{2d-1,2d\}=\{3,4\}$ driven by $\mathrm{sgn}(\Delta_2)$.

The companion conjecture introduced in \textsc{channel-theory-v11},
\emph{op:xi0-degree-d}, asks whether the proven $d=2$ identity
$\xi_0(b)=2/\sqrt{\beta_2}$ (Proposition~3.3.A) extends to all $d\ge 1$
in the form $\xi_0(b)=c(d)/\beta_d^{1/d}$.

\section{Catalogue}

We enumerate $b(n) = a_4 n^4 + a_3 n^3 + a_2 n^2 + a_1 n + a_0$ with
$a_4\in\{1,2,3,5,7\}$, $a_{0,1,2,3}\in\{-3,\ldots,3\}$ in lex order
($a_4, a_3, a_2, a_1, a_0$), keeping $\gcd=1$ and irreducible
over $\mathbb{Q}$, until $60$ entries. Galois group distribution:
$57\times S_4 + 3\times D_4 + 0\times V_4/C_4/A_4$. Real signature:
$(r_1,r_2)=(2,1)$ for $59$, $(4,0)$ for $1$. The catalogue contains
no CM-quartic ($r_1=0$). The window is dominated by mixed $S_4$;
this is a structural fact about the lex window and not a sampling
artefact.

Calibration anchors $x^4-2\ (D_4),\ x^4+1\ (V_4),\ x^4-x-1\ (S_4)$
all reproduce expected Galois group.

\section{Empirical evidence for Conjecture~B4 at $d=4$}

For each cataloged family we compute $L_b$ at $\mathrm{dps}=1200$
for $N\in[10,60]$ step $2$, with reference $N_{\mathrm{ref}}=150$, and
fit the ansatz of Conjecture~\ref{conj:B4-sharp} by linear LSQ in
$(\log N\cdot N,\,N,\,\log N,\,1)$. The protocol matches
Session~B's $d=3$ harvest with adjusted window $[10,60]$ (vs.\ $[10,100]$)
because at $d=4$ the $A\,N\log N$ decay is faster and $N=60$ already
yields $|\delta_N|\sim 10^{-700}$ within the working precision floor.

\paragraph{Results (Phase~Q1-2).}
\begin{itemize}
  \item Verdict counts: $\mathtt{b4\_consistent\_at\_8}: 60/60$;
        $\mathtt{b4\_split\_candidate\_at\_7}: 0$;
        $\mathtt{b4\_violating}: 0$.
  \item Empirical $A$-distribution: $\overline{A}=7.9544$,
        $\sigma_A=0.0037$, $A_{\min}=7.9488$, $A_{\max}=7.9605$.
  \item All $60$ families satisfy $|A_{\mathrm{fit}}-8|\le 0.0512 < 0.10$
        and $|A_{\mathrm{fit}}-7|\ge 0.95\gg 0.30$.
  \item Stderr range $[0.0021,\,0.0028]$ uniformly.
\end{itemize}

\begin{remark}[Systematic offset]\label{rem:offset}
The cluster $\overline{A}=7.954$ sits $0.046$ below $A=8$. The same
4-parameter ansatz at $d=3$ produced $\overline{A}=5.973$
($\sigma=0.026$, Session~B), at $d=2$ produced $\overline{A}=3.97$ on the
$A=4$ stratum (PCF-1 v1.3). The offset scales roughly as $0.05$ per
unit increase in $d$ and is consistent with a fit-window bias from
the absence of a $-(\log\log N)$ correction in the ansatz. None of
the $60$ families in this catalogue is thereby pushed toward $A=7$.
\end{remark}

\section{Universal $\xi_0$-identity at $d=4$}

For a degree-$d$ PCF with $a_n=1$ and
$b(n)=\sum_{k=0}^{d}\beta_k n^k$, the partial-denominator generating
function $f(z)=\sum_n Q_n z^n$ satisfies
\[
L_d\,f \;=\; (\text{polynomial in } z\text{ of degree }\le 1),\qquad
L_d \;=\; 1 - z\,B_d(\theta+1) - z^2,\quad \theta=z\partial_z,
\]
where $B_d(t)=\sum_{k=0}^{d}\beta_k\,t^k$.

\begin{proposition}[Universal $\xi_0$-identity, derived]\label{prop:xi0-univ}
The Newton polygon of $L_d$ at $z=0$ has a unique nontrivial slope-$(1/d)$
edge $E:(0,0)\to(1,d)$. The characteristic polynomial along $E$ under
the WKB ansatz $f\sim\exp(c/u)$, $z=u^d$, is
\[
\chi_d(c)\;=\;1\;-\;\frac{\beta_d}{d^d}\,c^d.
\]
Its leading positive real root is
\[
\boxed{\ \xi_0(b)\;=\;\frac{d}{\beta_d^{1/d}}\;=\;c(d)/\beta_d^{1/d}\ \text{with}\ c(d)=d\ }.
\]
\end{proposition}

\begin{proof}
$\theta\,\exp(c/u)=(u/d)\,\partial_u\exp(c/u)=-(c/(d u))\,\exp(c/u)$.
So $\theta^j$ acts as $(-c/(d u))^j$ at leading order. The single
edge $E$ of $L_d$ at slope $1/d$ joins $(0,0)$ (coefficient $+1$)
and $(1,d)$ (coefficient $-\beta_d$). Substituting $z=u^d$ and
collecting the $u^0$ coefficient of $L_d \exp(c/u)$ gives
$1 - \beta_d (-c)^d/d^d=0$, equivalently $1=\beta_d c^d/d^d$.
The positive real root is $c=d/\beta_d^{1/d}$.
\end{proof}

\paragraph{Numerical verification at $d=4$ (Phase~Q1-3).}
On $8$ representative quartics --- one per available bin
($\mathtt{mixed\_S_4}$ smallest $|\Delta_4|$, $\mathtt{mixed\_S_4}$
largest $|\Delta_4|$, $\mathtt{mixed\_D_4}$, $\mathtt{+\_real\_D_4}$),
three out-of-window anchors ($x^4-2,\ x^4+1,\ x^4-x-1$), and one
$a_4=7$ sample --- the inferred constant
$c(4):=\xi_0(b)\,a_4^{1/4}$ takes the value $4$ to spread $0$ at
$\mathrm{dps}=80$. The conjecture is verified at $d=4$ with
$c(4)=d=4$, as predicted by Proposition~\ref{prop:xi0-univ}.

\section{Cross-degree compilation and refined op:d2-anomaly}

\begin{table}[h]
\centering
\begin{tabular}{cllll}
\hline
$d$ & $A$ branches predicted & $A$ branches empirical & $\xi_0(b)$ & source \\
\hline
$2$ & $\{3,4\}$ & $\{3,4\}$ split by $\mathrm{sgn}(\Delta_2)$ &
$2/\sqrt{\beta_2}$ & PCF-1 v1.3 Thm 5 \\
$3$ & $\{5,6\}$ & $\{6\}$ alone (50/50) &
$3/\beta_3^{1/3}$ (predicted) & PCF-2 v1.1 \\
$4$ & $\{7,8\}$ & $\{8\}$ alone (60/60) &
$4/\beta_4^{1/4}$ (verified) & this note \\
\hline
\end{tabular}
\end{table}

\paragraph{Refined op:d2-anomaly status.}
The $\{2d-1,2d\}$ split at $d=2$ is so far the unique populated
multi-branch instance. Both $d=3$ (Session~B/C1) and $d=4$ (this note)
populate only the sharp branch $A=2d$ within their respective lex
windows.

\paragraph{Caveats.}
(i)~The $d=4$ catalogue contains no CM quartic ($r_1=0$, $V_4/C_4$);
a CM-quartic harvest is the natural follow-up.  (ii)~The
$\overline{A}\approx 2d-0.05$ cluster is fit-window bias, not a
genuine displacement (Remark~\ref{rem:offset}).  (iii)~Universality
of $\xi_0$ is an \emph{operational} (Newton-polygon) statement, not
yet a Borel-radius identity; the latter remains open per
\textsc{cc-pipeline-g} / \textsc{borel-channel-5x}.

\section{Conjecture B4 at $d\ge 3$: status after Q1}

Combining Sessions~B/C1 ($d=3$) and Q1 ($d=4$):
$50/50$ cubic $+\,60/60$ quartic $=\,110/110$ families verdict
\emph{$A=2d$, no $A=2d-1$ branch detected} within the joint lex
window. Conjecture~\ref{conj:B4-sharp} (B4 sharp form for $d\ge 3$)
remains the working program statement; B4$'$ (the conjectured $A=2d-1$
elementary-positive branch) is empirically empty in both degrees.

\subsection*{Acknowledgments.}
This is a Copilot-pipeline session output (SIARC). Strategic direction
and final review are by Claude (claude.ai); execution and code are by
GitHub Copilot. The cubic baseline is from PCF-2 v1.1 Sessions B and
C1; the $\xi_0$ proposition recovers and generalises Proposition~3.3.A
of CHANNEL-THEORY-V11.

\end{document}
